% =====================================================================
% 別冊『Git & GitHub 応用 — Issue と自動化でチーム開発を回す』
% エンジン: LuaLaTeX（jlreq / book）
% ビルド:   latexmk main_advanced.tex   （出力は output/main_advanced.pdf）
%           全冊まとめて作るときは、このフォルダで latexmk のみ実行
%
% 本文のスクリーンショットは \shotplaceholder で仮置き（images/04_advanced/ に置いて差し替える）。
% 対象: 別冊『チームで安全に開発する』を終え、運用を一段引き上げたい人。
% =====================================================================
\documentclass[book, paper=a4, fontsize=10.5pt, titlepage]{jlreq}

% =====================================================================
% preamble.tex — パッケージ読込・色・カスタム環境・マクロの定義
% 本文（chapters/*.tex）から使う見た目の部品はすべてここに集約する。
% エンジン: LuaLaTeX + jlreq
% =====================================================================

% --- 基本パッケージ ---------------------------------------------------
\usepackage{graphicx}          % 画像の取り込み
\usepackage{xcolor}            % 色指定
\usepackage{booktabs}          % 罫線のきれいな表
\renewcommand{\arraystretch}{1.25} % 表の行を少しゆったりさせる
\usepackage{tabularx}          % 幅を自動調整する表
\usepackage{enumitem}          % 箇条書きの調整
\usepackage{fancyvrb}          % コード枠に使う Verbatim
\usepackage[most]{tcolorbox}   % 装飾枠（コード・注意書き）
\usepackage{adjustbox}         % 図が本文幅を超えたときの自動縮小
\usepackage{tikz}              % 概念図・フロー図
\usetikzlibrary{arrows.meta,positioning,shapes.geometric,calc,fit,backgrounds}
\usepackage{float}             % 図を本文の位置に固定する [H]
\usepackage{caption}           % 図番号つきキャプションの体裁
\captionsetup{font=small, labelfont=bf, labelsep=quad}

% 画像の探索パス（既存の Markdown 版と同じ images/ を共有する）
\graphicspath{{../images/}{./images/}}

% --- 配色 -------------------------------------------------------------
% Markdown 版の Mermaid 図の色をそのまま踏襲する（薄い色は使わない方針）。
\definecolor{gitRed}{HTML}{C0392B}     % 作業ツリー・注意
\definecolor{gitOrange}{HTML}{D68910}  % ステージング
\definecolor{gitGreen}{HTML}{2E7D32}   % リポジトリ・完了
\definecolor{gitBlue}{HTML}{4A90D9}    % コミット・情報
\definecolor{gitPurple}{HTML}{6C3483}  % 自分側
\definecolor{gitTeal}{HTML}{117A65}    % 両方採用
\definecolor{gitAmber}{HTML}{8A6D0B}   % 判断ノード
\definecolor{codeBg}{HTML}{F4F5F7}     % コード枠の背景
\definecolor{codeFrame}{HTML}{D0D3D9}  % コード枠の枠線

% --- コード枠 ---------------------------------------------------------
% 使い方: \begin{codebox}[bash] ... Verbatim ... \end{codebox}
% 省略時のラベルは bash。右上に言語ラベルを小さく表示する。
\newtcolorbox{codebox}[1][bash]{%
  enhanced, breakable,
  colback=codeBg, colframe=codeFrame, boxrule=0.6pt, arc=2pt,
  left=9pt, right=9pt, top=13pt, bottom=8pt,
  before skip=8pt, after skip=8pt,
  overlay={%
    \node[anchor=north east, font=\sffamily\scriptsize\bfseries,
          text=gray!60] at ([xshift=-6pt, yshift=-3pt]frame.north east) {#1};%
  },
}
% コード本体を書くための Verbatim 環境（小さめ・行間やや詰め）
\DefineVerbatimEnvironment{code}{Verbatim}{fontsize=\small, xleftmargin=0pt}

% 本文中のインラインコード: \code{git status}
\newcommand{\cmd}[1]{\texttt{#1}}

% --- 注意書き（admonition）枠 ----------------------------------------
% ヒント（緑）・注意（赤）・補足（グレー）の3種。emoji は使わずラベルで表す。
% タイトルは本文（明るい背景）に直接乗るため、白文字だと見えない。
% タイトル文字色は各テーマ色にして、はっきり読めるようにする。
\newtcolorbox{hintbox}{%
  enhanced, breakable,
  colback=gitGreen!5, colframe=gitGreen, boxrule=0.8pt, arc=2pt,
  left=8pt, right=8pt, top=6pt, bottom=6pt,
  coltitle=gitGreen, fonttitle=\sffamily\bfseries\small,
  title=ヒント, attach title to upper=\quad,
}
\newtcolorbox{warnbox}{%
  enhanced, breakable,
  colback=gitRed!5, colframe=gitRed, boxrule=0.8pt, arc=2pt,
  left=8pt, right=8pt, top=6pt, bottom=6pt,
  coltitle=gitRed, fonttitle=\sffamily\bfseries\small,
  title=注意, attach title to upper=\quad,
}
% 執筆メモ（雛形用の枠）。原稿を清書したら、この枠ごと削除する。
\newtcolorbox{draftbox}{%
  enhanced, breakable,
  colback=gitAmber!6, colframe=white, boxrule=0pt, arc=2pt,
  borderline={0.9pt}{0pt}{gitAmber, dashed},
  left=8pt, right=8pt, top=6pt, bottom=6pt,
  coltitle=gitAmber, fonttitle=\sffamily\bfseries\small,
  title=執筆メモ, attach title to upper=\quad,
}
\newtcolorbox{notebox}{%
  enhanced, breakable,
  colback=gray!8, colframe=gray!60!black, boxrule=0.8pt, arc=2pt,
  left=8pt, right=8pt, top=6pt, bottom=6pt,
  coltitle=gray!30!black, fonttitle=\sffamily\bfseries\small,
  title=補足, attach title to upper=\quad,
}

% --- スクリーンショット取り込み --------------------------------------
% \screenshot[幅の割合]{ファイル}{キャプション}
% フロートにせず本文の流れに沿って中央寄せで置く（説明の直後に見せたいため）。
\newcommand{\screenshot}[3][0.7]{%
  \begin{center}
    \includegraphics[width=#1\linewidth]{#2}\\[3pt]
    {\small\itshape #3}%
  \end{center}%
}

% \screenshottrim[幅の割合]{ファイル}{キャプション}{trim指定}
% 余白の多いスクリーンショットを切り詰めて大きく見せる。
% trim は「左 下 右 上」の順で、画像1ピクセル = 1bp として指定する
% （例: 1920x1080 の下半分を捨てる → {0 540bp 0 0}）。
\newcommand{\screenshottrim}[4][0.95]{%
  \begin{center}
    \includegraphics[width=#1\linewidth, trim=#4, clip]{#2}\\[3pt]
    {\small\itshape #3}%
  \end{center}%
}

% --- TikZ 共通スタイル ------------------------------------------------
% Mermaid のフロー図を再現するためのノード／矢印スタイル。
\tikzset{
  flownode/.style={%
    rounded corners=4pt, draw=none, text=white,
    font=\sffamily\small, align=center,
    minimum height=1.15cm, minimum width=2.2cm, inner sep=6pt},
  decisionnode/.style={%
    diamond, aspect=1.9, draw=none, text=white,
    font=\sffamily\small, align=center, inner sep=3pt},
  flowarrow/.style={-{Latex[length=3.6mm, width=3mm]}, line width=1.8pt, draw=black},
  edgelabel/.style={font=\sffamily\normalsize\bfseries, fill=white, inner sep=3pt},
}
% 図全体を中央寄せで置くためのヘルパ環境（図番号なし）。
% 本文幅を超えた図は自動で縮小する（幅が足りている図は拡大しない）。
\newenvironment{diagram}%
  {\par\medskip\begin{center}\begin{adjustbox}{max width=\linewidth}}%
  {\end{adjustbox}\end{center}\par\medskip}

% --- 図番号つきの図（本文が「図1.2のように」と参照する概念図・フロー図用）---
% 手順に紐づく操作スクショには使わない（そちらは \screenshot を使う）。
\makeatletter
% TikZ など本体を囲む環境: \begin{numfig}{キャプション}{ラベル} ... \end{numfig}
\newenvironment{numfig}[2]{%
  \def\nf@cap{#1}\def\nf@lab{#2}%
  \begin{figure}[H]\centering\begin{adjustbox}{max width=\linewidth}%
}{%
  \end{adjustbox}\caption{\nf@cap}\label{\nf@lab}\end{figure}%
}
\makeatother
% 画像1枚を図番号つきで置く: \figimage[幅]{ファイル}{キャプション}{ラベル}
\newcommand{\figimage}[4][0.95]{%
  \begin{figure}[H]\centering
    \includegraphics[width=#1\linewidth]{#2}
    \caption{#3}\label{#4}
  \end{figure}%
}


% hyperref はプリアンブルの最後（他パッケージ読込後）に置くのが安全。
\usepackage[unicode, hidelinks]{hyperref}

% --- 書誌情報 ---------------------------------------------------------
\title{Git \& GitHub 応用}
\author{}
\date{\today}

\begin{document}

% --- 表紙 -------------------------------------------------------------
\begin{titlepage}
  \centering
  \vspace*{4cm}
  {\sffamily\bfseries\Huge Git \& GitHub 応用\par}
  \vspace{1.2cm}
  {\sffamily\Large Issue と自動化でチーム開発を回す\par}
  \vspace{0.6cm}
  {\sffamily\normalsize ── 別冊『チームで安全に開発する』の続き ──\par}
  \vfill
  {\large \today\par}
\end{titlepage}

% --- この本のねらい ---------------------------------------------------
\thispagestyle{empty}
\vspace*{1.5cm}
{\sffamily\LARGE\bfseries はじめに —この本を読むとできること—\par}
\vspace{1em}

本書は、別冊『チームで安全に開発する』の\textbf{続き}です。ブランチとPRでの共同作業ができるようになった人が、\textbf{チームの進め方そのものを整える}ための冊子です。読み終えると、次のことができるようになります。

\begin{itemize}[leftmargin=1.6em]
  \item \textbf{Issue} で「やること」を見える化し、担当と期限を管理できる
  \item \textbf{マイルストーン}で締め切りを決め、完了率で遅れに早く気づける
  \item \textbf{Projects} で、Issueをボードや\textbf{ガントチャート}にして進み具合を共有できる
  \item Issue とブランチ・PR を結びつけ、\textbf{マージと同時にIssueを閉じられる}
  \item \textbf{GitHub Actions} で、\cmd{push} やPRのたびに自動チェックを走らせられる
  \item ブランチ保護と組み合わせ、\textbf{チェックが通らないとマージできない}状態にできる
\end{itemize}

つまり、\textbf{「人が気をつける」から「仕組みで守る」へ}進むところまでが、本書のゴールです。

\begin{warnbox}
\textbf{前提：} 本書は、別冊『チームで安全に開発する』の内容（ブランチ・PR・コンフリクト解決）が済んでいる前提で進みます。まだの人は、先にそちらを読んでください。
\end{warnbox}

\begin{hintbox}
\textbf{この本の読み方}
\begin{itemize}[leftmargin=1.4em, nosep]
  \item 作業の管理（Issue・マイルストーン・ガントチャート）を整えたい人 → 第1章へ
  \item PRのたびに自動でテストを走らせたい人 → 第2章へ（第1章の「IssueとPRをつなぐ」だけ先に読むと流れがつかみやすい）
  \item 本書の先に何を学べばよいか知りたい人 → 付録A「この先に進むための地図」へ
\end{itemize}
\end{hintbox}

\clearpage

% --- 目次 -------------------------------------------------------------
\tableofcontents
\clearpage

% --- 本文 -------------------------------------------------------------
% =====================================================================
% 応用編 第1章 Issueで作業を管理する
% ステータス: 雛形（見出しと執筆メモのみ。本文はこれから書く）
% =====================================================================
\chapter{Issueで作業を管理する}

別冊『チームで安全に開発する』では、\textbf{ブランチを分けてPRでマージする}という「変更の流し方」を扱いました。この章で扱うのは、その1つ手前にある\textbf{「何を、誰が、いつやるのか」}の管理です。GitHubでは、これを\textbf{Issue（イシュー）}という機能で行います。

\begin{draftbox}
\textbf{この章のゴール}（清書時にこの枠は削除する）
\begin{itemize}[leftmargin=1.4em, nosep]
  \item Issueを立てて、作業をチームで見える化できる
  \item ラベル・担当者・マイルストーンで、Issueを整理できる
  \item IssueとPRを結びつけ、マージと同時にIssueを閉じられる
  \item Projects（ボード）で進捗を一覧できる
\end{itemize}
\textbf{想定読者}：別冊『チームで安全に開発する』を終えた人。ブランチとPRは分かっている前提で書く。
\end{draftbox}

\begin{hintbox}
\textbf{この章の読み方}
\begin{itemize}[leftmargin=1.4em, nosep]
  \item まず動かしたい人 → 「Issueを立てる」→「IssueとPRをつなぐ」だけ読めばOK
  \item チームの運用ルールを決めたい人 → 「運用の型」へ
\end{itemize}
\end{hintbox}

% ---------------------------------------------------------------------
\section{Issueとは何か}

\begin{draftbox}
\textbf{書くこと}
\begin{itemize}[leftmargin=1.4em, nosep]
  \item Issue＝\textbf{リポジトリに貼る「作業の付箋」}。バグ報告・機能要望・やることメモを1件1枚で管理する
  \item PRとの違い：\textbf{Issueは「やること」、PRは「やった変更」}。Issueが先、PRが後
  \item チャットやメールで作業を頼むのとの違い：\textbf{リポジトリに残り、検索でき、変更履歴とひも付く}
  \item シリーズ共通の倉庫のたとえに接続する：Issueは「\textbf{倉庫に貼る作業指示書}」
\end{itemize}
\textbf{入れる図}：Issue（やること）→ ブランチ → PR（やった変更）→ マージ、の流れ図。
TikZ の \cmd{flownode} / \cmd{flowarrow} を使い、既存の図と体裁を揃える。
\end{draftbox}

% ---------------------------------------------------------------------
\section{Issueを立てる}

\begin{draftbox}
\textbf{書くこと}：GitHubのIssuesタブ →「New issue」→ タイトルと本文を書いて作成、までの手順。
\textbf{必要なスクリーンショット}（\cmd{images/04\_advanced/} に置く）
\begin{itemize}[leftmargin=1.4em, nosep]
  \item \cmd{issue\_new\_button.png}：Issuesタブと「New issue」ボタン
  \item \cmd{issue\_form.png}：タイトル・本文の入力画面
  \item \cmd{issue\_created.png}：作成後のIssue画面（番号 \cmd{\#12} が付くことを示す）
\end{itemize}
\end{draftbox}

\subsection{タイトルと本文の書き方}

\begin{draftbox}
\textbf{書くこと}
\begin{itemize}[leftmargin=1.4em, nosep]
  \item タイトルは\textbf{「何をどうしたいか」が一言で分かる形}にする（例：\cmd{ログイン画面のエラー表示が出ない}）
  \item 本文に書く3点：\textbf{現状／期待する状態／再現手順（またはやること）}
  \item 悪い例と良い例を並べて示す（「バグ直す」→ 何のバグか分からない）
  \item コミットメッセージの書き方（本編）と同じ考え方であることに触れる
\end{itemize}
\end{draftbox}

\subsection{Issueテンプレートを用意する}

\begin{draftbox}
\textbf{書くこと}：\cmd{.github/ISSUE\_TEMPLATE/} にテンプレートを置くと、New issue のときに雛形が出る。
付録A（別冊のPRテンプレート）と同じ形式で、\textbf{そのままコピーして使えるテンプレート}を載せる。

最低限、次の2種類を用意する：
\begin{itemize}[leftmargin=1.4em, nosep]
  \item \cmd{bug\_report.md}：不具合の報告用（現象／再現手順／期待する動作／環境）
  \item \cmd{task.md}：やること用（背景／完了条件／関連Issue）
\end{itemize}
\end{draftbox}

\begin{codebox}[markdown]
\begin{code}
（ここに bug_report.md の中身を貼る）
---
name: バグ報告
about: 動かない・想定と違う挙動を報告する
labels: bug
---

## 現象

## 再現手順

## 期待する動作

## 環境
\end{code}
\end{codebox}

% ---------------------------------------------------------------------
\section{ラベル・担当者・マイルストーンで整理する}

\begin{draftbox}
\textbf{書くこと}：Issueが増えたときに迷子にならないための3つの道具。

\begin{itemize}[leftmargin=1.4em, nosep]
  \item \textbf{ラベル}：種類で色分けする（\cmd{bug} / \cmd{enhancement} / \cmd{question} / \cmd{good first issue}）。おすすめの初期セットを表で示す
  \item \textbf{担当者（Assignees）}：誰がやるかを1人に決める。空のままにしない
  \item \textbf{マイルストーン}：締め切りやリリース単位でまとめる（例：\cmd{v1.0 提出}）
\end{itemize}
\textbf{入れる表}：ラベル名・色・使う場面の3列。
\textbf{必要なスクリーンショット}：\cmd{issue\_labels.png}（サイドバーでラベル・担当者を設定している画面）。
\end{draftbox}

% ---------------------------------------------------------------------
\section{IssueとPRをつなぐ}

\begin{draftbox}
\textbf{この節がこの章の山場}。PR本文に \cmd{Closes \#12} と書くと、\textbf{マージと同時にIssueが自動で閉じる}ことを示す。

\textbf{書くこと}
\begin{itemize}[leftmargin=1.4em, nosep]
  \item 使えるキーワード（\cmd{Closes} / \cmd{Fixes} / \cmd{Resolves}）と、単に参照だけしたいときの \cmd{\#12}
  \item ブランチ名にIssue番号を入れる運用（例：\cmd{feature/12-add-login}）
  \item Issue画面から直接ブランチを作る方法（\textbf{Create a branch} リンク）
  \item 「Issueを立てる → ブランチを切る → PRを出す → マージでIssueが閉じる」の一周を、実際の画面で通す
\end{itemize}
\textbf{入れる図}：この一周のフロー図（別冊 第1章のPRフロー図と体裁を揃える）。
\textbf{必要なスクリーンショット}：\cmd{issue\_linked\_pr.png}（PR本文の \cmd{Closes \#12} と、Issue側に出るリンク表示）。
\end{draftbox}

% ---------------------------------------------------------------------
\section{Projectsでボードにする}

\begin{draftbox}
\textbf{書くこと}：Issueを「未着手／作業中／レビュー中／完了」の列に並べて、進捗を一覧する。

\begin{itemize}[leftmargin=1.4em, nosep]
  \item Projects（ボード）の作り方と、Issueの追加方法
  \item ステータスの自動更新（PRがマージされたら「完了」に動かす設定）
  \item 少人数なら\textbf{ラベル運用だけでも足りる}こと。Projectsは人数と件数が増えてから入れるのがおすすめ、と\textbf{使いどころを限定}して書く
\end{itemize}
\textbf{必要なスクリーンショット}：\cmd{project\_board.png}（ボード全体）。
\end{draftbox}

% ---------------------------------------------------------------------
\section{運用の型}

\begin{draftbox}
\textbf{書くこと}：そのまま真似できる運用ルールを、別冊 付録Aと同じ調子で示す。

\begin{itemize}[leftmargin=1.4em, nosep]
  \item \textbf{作業を始める前に必ずIssueを立てる}（口頭で決めたことも1行で残す）
  \item \textbf{1 Issue = 1 ブランチ = 1 PR} を基本にする
  \item 週に1度、\textbf{開いたままのIssueを棚卸し}する（閉じ忘れの防止）
  \item 議論はチャットではなく\textbf{Issueのコメントに書く}（後から経緯を追えるようにする）
\end{itemize}
\end{draftbox}

% ---------------------------------------------------------------------
\section*{この章のまとめ}
\addcontentsline{toc}{section}{この章のまとめ}

\begin{draftbox}
本文を書き終えてから、次の骨子でまとめを作る。
\begin{itemize}[leftmargin=1.4em, nosep]
  \item Issueは「やること」、PRは「やった変更」
  \item Issueは\textbf{タイトルで内容が分かる}ように書く
  \item \cmd{Closes \#12} でPRとIssueをつなぐと、マージで自動的に閉じる
  \item 1 Issue = 1 ブランチ = 1 PR
\end{itemize}
次章は、この流れに\textbf{自動チェック}を足す話（GitHub Actions）につなげる。
\end{draftbox}

% =====================================================================
% 応用編 第2章 GitHub Actionsで自動化する
% ステータス: 雛形（見出しと執筆メモのみ。本文はこれから書く）
% =====================================================================
\chapter{GitHub Actionsで自動化する}

前章までで、\textbf{Issueで作業を決め、ブランチを切り、PRでマージする}という流れができました。この章では、その流れの上に\textbf{自動チェック}を載せます。「\cmd{push} するたびにテストが走る」「PRを出すと、通らない限りマージできない」という状態を、\textbf{GitHub Actions} で作ります。

\begin{draftbox}
\textbf{この章のゴール}（清書時にこの枠は削除する）
\begin{itemize}[leftmargin=1.4em, nosep]
  \item ワークフローファイル（\cmd{.github/workflows/*.yml}）の構造を読める
  \item \cmd{push} / PR をきっかけに動く、最初のワークフローを作れる
  \item 失敗したときに、どこを見て直せばよいか分かる
  \item ブランチ保護と組み合わせて、\textbf{チェックが通らないとマージできない}状態にできる
\end{itemize}
\textbf{書くときの注意}：Actionsは機能が広いので、\textbf{「PRのたびに自動チェックを走らせる」1本に絞る}。
デプロイ・リリース自動化には深入りせず、最後に「次に何を学べばよいか」として名前だけ挙げる。
\end{draftbox}

\begin{hintbox}
\textbf{この章の読み方}
\begin{itemize}[leftmargin=1.4em, nosep]
  \item とりあえず動かしたい人 → 「最初のワークフローを作る」を写経すればOK
  \item YAMLの読み方が不安な人 → 先に「ワークフローファイルの構造」へ
\end{itemize}
\end{hintbox}

% ---------------------------------------------------------------------
\section{GitHub Actionsとは何か}

\begin{draftbox}
\textbf{書くこと}
\begin{itemize}[leftmargin=1.4em, nosep]
  \item Actions＝\textbf{GitHub上に住んでいる「自動作業係」}。「\cmd{push} されたら」「PRが出たら」といった\textbf{きっかけ}で、決めておいた作業を実行する
  \item 動く場所は\textbf{GitHubが用意する使い捨てのPC（ランナー）}であり、自分のPCではないこと
  \item うれしいこと：\textbf{レビュー前に機械が確認してくれる}／\textbf{人によって手順がぶれない}／\textbf{確認忘れがなくなる}
  \item シリーズのたとえに接続：Actionsは「\textbf{倉庫に荷物を入れる前の検品係}」
  \item 料金の注意：Publicリポジトリは無料、Privateは無料枠あり（\textbf{枠を超えると課金}）。最新の条件は公式を見るよう促す
\end{itemize}
\textbf{入れる図}：push / PR（きっかけ）→ ランナー起動 → ジョブ実行 → 結果がPRに緑✓／赤×で出る、の流れ図。
\end{draftbox}

% ---------------------------------------------------------------------
\section{ワークフローファイルの構造}

\begin{draftbox}
\textbf{書くこと}：\cmd{.github/workflows/} に置いたYAMLファイルが、そのままワークフローになる。
下の骨組みを1行ずつ分解して説明する（\cmd{name} / \cmd{on} / \cmd{jobs} / \cmd{runs-on} / \cmd{steps} / \cmd{uses} / \cmd{run}）。

\textbf{入れる表}：キーワードと意味の対応表（\cmd{on}＝いつ動くか、\cmd{jobs}＝何をするかのまとまり、
\cmd{steps}＝ジョブの中の手順、\cmd{uses}＝公開部品を使う、\cmd{run}＝コマンドを実行する）。

\textbf{YAMLの注意}：\textbf{字下げが意味を持つ}こと、タブではなく半角スペースを使うことを警告枠で。
\end{draftbox}

\begin{codebox}[yaml]
\begin{code}
# .github/workflows/ci.yml
name: CI                      # ワークフローの名前（Actionsタブに表示される）

on:                           # いつ動かすか（きっかけ）
  push:
    branches: [ main ]
  pull_request:

jobs:                         # 何をするか
  test:                       # ジョブの名前（自由に付けられる）
    runs-on: ubuntu-latest    # どのOSのランナーで動かすか
    steps:                    # 上から順に実行される手順
      - uses: actions/checkout@v4   # リポジトリを取ってくる（定番の1手目）
      - run: echo "ここに実際のチェックを書く"
\end{code}
\end{codebox}

% ---------------------------------------------------------------------
\section{最初のワークフローを作る}

\begin{draftbox}
\textbf{書くこと}：写経して1本動かすまで。ここが章の中心。

\textbf{手順の骨子}
\begin{enumerate}[leftmargin=1.6em, nosep]
  \item ブランチを切る（\cmd{feature/add-ci}）
  \item \cmd{.github/workflows/ci.yml} を作る（中身は上の骨組み）
  \item \cmd{add} → \cmd{commit} → \cmd{push} → PRを出す
  \item PR画面の下部に\textbf{チェック結果}が出るのを確認する
  \item 緑✓になったらマージする
\end{enumerate}

\textbf{必要なスクリーンショット}（\cmd{images/04\_advanced/} に置く）
\begin{itemize}[leftmargin=1.4em, nosep]
  \item \cmd{actions\_tab.png}：Actionsタブの実行一覧
  \item \cmd{pr\_checks\_running.png}：PR下部で実行中のチェック
  \item \cmd{pr\_checks\_passed.png}：緑✓になった状態
\end{itemize}

\textbf{題材の候補}：本シリーズの読者に合わせるなら、\textbf{LaTeX原稿のビルド}か\textbf{textlintによる文章チェック}が題材として自然。
Pythonのテスト（\cmd{pytest}）を例にするかは、想定読者を決めてから選ぶ。
\end{draftbox}

% ---------------------------------------------------------------------
\section{よくある使いどころ}

\begin{draftbox}
\textbf{書くこと}：「何を自動化すると効くか」を、\textbf{表で一覧}にする（きっかけ／やること／うれしさ の3列）。

候補：
\begin{itemize}[leftmargin=1.4em, nosep]
  \item PRのたびに\textbf{テスト}を走らせる
  \item PRのたびに\textbf{文法・書式チェック}（lint、textlint）を走らせる
  \item \cmd{main} が更新されたら\textbf{ドキュメント（PDF）をビルド}して成果物を残す
  \item タグを打ったら\textbf{リリースを作る}
\end{itemize}
それぞれ\textbf{どのきっかけ（\cmd{on}）で動かすか}を対応づけて書く。
\end{draftbox}

% ---------------------------------------------------------------------
\section{実行結果の見方と、失敗したときの直し方}

\begin{draftbox}
\textbf{書くこと}：赤×が出たときに読者が自力で進めるようにする節。

\begin{itemize}[leftmargin=1.4em, nosep]
  \item Actionsタブ → 対象の実行 → ジョブ → \textbf{失敗したステップを開いてログを読む}
  \item ログは\textbf{上から読むのではなく、赤い行と最後の数十行}を見るのがコツ
  \item よくある失敗を表に：YAMLの字下げミス／ファイルの置き場所が違う（\cmd{.github/workflows/} 直下か）／
        ブランチ名が \cmd{on} の条件と一致していない／依存パッケージの入れ忘れ
  \item 直したら\textbf{同じブランチに \cmd{push} すれば自動で再実行}される（PRを作り直さなくてよい）
\end{itemize}
\textbf{必要なスクリーンショット}：\cmd{actions\_failed\_log.png}（失敗ログを開いた画面）。
\end{draftbox}

% ---------------------------------------------------------------------
\section{ブランチ保護と組み合わせる}

\begin{draftbox}
\textbf{書くこと}：自動チェックを\textbf{マージの条件}にして、はじめて効果が出る。

\begin{itemize}[leftmargin=1.4em, nosep]
  \item Settings → Branches → \textbf{ブランチ保護ルール}で \cmd{main} を守る
  \item 設定する項目：\textbf{PR必須}／\textbf{チェックの成功を必須}／\textbf{レビュー承認を必須}／\cmd{main} への直接pushを禁止
  \item 別冊で書いた「\cmd{main} は直接編集しない」という\textbf{約束を、仕組みとして強制する}のがこの節、という位置づけを明示する
  \item 少人数のうちは厳しくしすぎない、という運用上の注意も書く
\end{itemize}
\textbf{必要なスクリーンショット}：\cmd{branch\_protection.png}（保護ルールの設定画面）。
\end{draftbox}

% ---------------------------------------------------------------------
\section{秘密情報の扱い（Secrets）}

\begin{draftbox}
\textbf{書くこと}：APIキーやトークンをワークフローで使うときの置き場所。

\begin{itemize}[leftmargin=1.4em, nosep]
  \item \textbf{YAMLに直接書かない}。Settings → Secrets and variables → Actions に登録し、\cmd{\$\{\{ secrets.NAME \}\}} で参照する
  \item ログに出さない工夫（\cmd{echo} しない）
  \item 本編『はじめてのバージョン管理』の \cmd{.gitignore} の話とつなげる：\textbf{秘密情報はリポジトリに置かない}という原則は同じ
\end{itemize}
\textbf{枠は warnbox を使う}（ヒントではなく注意）。
\end{draftbox}

% ---------------------------------------------------------------------
\section*{この章のまとめ}
\addcontentsline{toc}{section}{この章のまとめ}

\begin{draftbox}
本文を書き終えてから、次の骨子でまとめを作る。
\begin{itemize}[leftmargin=1.4em, nosep]
  \item Actionsは「きっかけ（\cmd{on}）」と「やること（\cmd{jobs}）」を書くだけ
  \item ファイルの置き場所は \cmd{.github/workflows/} 直下
  \item 失敗したら\textbf{赤いステップのログ}を読む。直して \cmd{push} すれば再実行される
  \item ブランチ保護と組み合わせて、はじめて「壊れたものがマージされない」状態になる
\end{itemize}
\textbf{次に学ぶとよいもの}として名前だけ挙げる：自作アクション、再利用可能ワークフロー、
マトリクスビルド、デプロイ（GitHub Pages など）。
\end{draftbox}


% --- 付録 -------------------------------------------------------------
\appendix

\chapter{この先に進むための地図}

本書では、範囲を\textbf{Issueによる作業の管理}と\textbf{PRのたびの自動チェック}に絞りました。この付録では、あえて扱わなかったものを\textbf{名前と一言だけ}並べます。手順は書きません。「こういうことがしたくなった」と思ったら、右の列の言葉でGitHubの公式ドキュメント（\cmd{docs.github.com}）を検索してください。

\section{作業の管理をもっと便利にする}

\begin{center}
\begin{tabularx}{\linewidth}{@{}>{\raggedright\arraybackslash}p{3.8cm} X >{\raggedright\arraybackslash}p{3.8cm}@{}}
\toprule
名前 & こんなときに & 調べる言葉 \\
\midrule
Issueフォーム & テンプレートを、入力欄やプルダウンのある\textbf{記入フォーム}にしたい & issue forms \\
\addlinespace[3pt]
Projectsのインサイト & 完了したIssueの数の推移を\textbf{グラフ}で見たい & project insights \\
\addlinespace[3pt]
GitHub CLI（\cmd{gh}） & Issueの作成やPRの確認を\textbf{ターミナルから}行いたい & gh issue, gh pr \\
\addlinespace[3pt]
CODEOWNERS & ファイルやフォルダごとに\textbf{レビュー担当を自動で割り当て}たい & CODEOWNERS \\
\bottomrule
\end{tabularx}
\end{center}

\section{自動化をもっと広げる}

\begin{center}
\begin{tabularx}{\linewidth}{@{}>{\raggedright\arraybackslash}p{3.8cm} X >{\raggedright\arraybackslash}p{3.8cm}@{}}
\toprule
名前 & こんなときに & 調べる言葉 \\
\midrule
デプロイの自動化 & \cmd{main} が更新されたら、Webページや文書を\textbf{自動で公開}したい & GitHub Pages, deployment \\
\addlinespace[3pt]
マトリクスビルド & 複数のOS・複数のPythonのバージョンで\textbf{まとめて}テストしたい & matrix strategy \\
\addlinespace[3pt]
キャッシュ & パッケージのインストールに時間がかかるので\textbf{速く}したい & dependency caching \\
\addlinespace[3pt]
再利用可能ワークフロー & 複数のリポジトリで\textbf{同じワークフローを使い回し}たい & reusable workflows \\
\addlinespace[3pt]
Dependabot & 使っているパッケージの更新を\textbf{自動でPRにして}ほしい & Dependabot \\
\addlinespace[3pt]
self-hosted runner & GPUのある研究室のPCなど、\textbf{自前のマシン}でActionsを動かしたい & self-hosted runners \\
\bottomrule
\end{tabularx}
\end{center}

\medskip

どれも、本書で身につけた\textbf{「Issueで決めて、PRで出して、機械に確かめさせる」}という流れの上に足していくものです。困りごとが出てきたときに、1つずつ試してみてください。

\end{document}
